\clearpage\thispagestyle{empty}
\pdfbookmark[0]{Sources for Gauss’s October 1804 Notice on Juno}{nachrichten1804-sources}
\begin{center}\textbf{Sources for Gauss’s October 1804 Notice on Juno}\end{center}
\noindent This AI translation follows Carl Friedrich Gauss’s notice as printed in \href{https://gdz.sub.uni-goettingen.de/id/PPN236007467}{\textit{Werke}, vol.~VI} (Göttingen, 1874), pp.~253--256. The \href{https://digipress.digitale-sammlungen.de/view/bsb10531607_00329_u001}{\textit{Braunschweigisches Magazin}}, no.~40, 6 October 1804, names Gauss in the byline in \href{https://www.digitale-sammlungen.de/en/view/bsb10531607?page=329}{col.~621}. The first printing ends with an ephemeris in \href{https://www.digitale-sammlungen.de/en/view/bsb10531607?page=332}{col.~628}. The collected edition refers readers to the \hyperlink{native-258}{ephemeris in Franz Xaver von Zach’s report on p.~253}.
\par\smallskip\noindent Karl Ludwig Harding discovered Juno. The star catalogue used by Harding was Jérôme de Lalande’s \href{https://archive.org/download/histoireclestef00lalagoog/page/n12.jpg}{\textit{Histoire céleste française}} (Paris, 1801). Gauss names Michel Lefrançais Lalande for his observations and John Flamsteed for his star catalogue.
\par\smallskip\noindent Translation and typesetting produced by AI. Magazine scans from the Bayerische Staatsbibliothek.
\clearpage
